\documentclass[12pt]{amsart}

\usepackage{amsmath}
\usepackage{amssymb}
\usepackage{amsthm}
\usepackage[dvipsnames]{xcolor}
\usepackage{hyperref}
\usepackage{booktabs}
\usepackage{array}

\hypersetup{
  colorlinks=true,
  linkcolor=MidnightBlue,
  citecolor=MidnightBlue,
  urlcolor=MidnightBlue
}

\newtheorem{theorem}{Theorem}[section]
\newtheorem{conjecture}[theorem]{Conjecture}
\theoremstyle{definition}
\newtheorem{definition}[theorem]{Definition}
\newtheorem{remark}[theorem]{Remark}

\title[CM Predicate for Degree-2 PCFs]{Complex Multiplication as a Transcendence Predicate
for Degree-2 Polynomial Continued Fractions}

\author{Papanokechi}
\thanks{ORCID: \href{https://orcid.org/0009-0000-6192-8273}{0009-0000-6192-8273}.
\textbf{Version 1.3 (2026-05-01):} folds in the WKB exponent identity
(Theorem~\ref{thm:wkb}, Sessions~E and~E$'$) and the channel-scope
subsection for the V$\_$quad prototype (Sessions~D, E, E$'$);
Conjecture~A part~(iv) is restated per the post-E$'$ decision matrix
(Variant~A: both $K=12$ marginals, QL15 and QL26, confirmed as
Pad\'e--Borel artefacts).
\textbf{Version 1.2 (2026-05-01):} added Session~C
(STOKES-FAMILY-EXTEND-3X) results for QL02, QL06, QL26;
Stokes predicate now confirmed across all six known
$\Delta<0$ degree-2 PCF families.
\textbf{Version 1.1 (2026-05-01):} added Session~B results for family
QL15 ($\Delta=-20$, non-Heegner). Version~1.0 published as
\href{https://doi.org/10.5281/zenodo.19931636}{DOI 10.5281/zenodo.19931636}.
\textbf{AI use disclosure.} This research was conducted using the
AEAL/SIARC multi-agent methodology. GitHub Copilot (powered by Anthropic
Claude Opus 4.7) and Anthropic Claude were used to assist with code
generation for numerical computations, cross-checking of intermediate
results, and drafting of expository prose. Sessions~D, E, and E$'$ of
the AEAL chain were executed via Anthropic Claude and GitHub Copilot
agents. All conjecture formulations, mathematical interpretations,
evidence-class assignments, and final editorial decisions were made by
the human author, who takes full responsibility for the content of this
manuscript. The AI tools are not authors. Session-level computation
logs and prompts are archived in the SIARC research workspace for
reproducibility per the AEAL methodology of~\cite{P8AEAL}.}
\address{Independent Researcher, Yokohama, Japan}
\email{}

\date{May 1, 2026}

\subjclass[2020]{Primary 11J70, 11J81; Secondary 34M55, 11G15}
\keywords{polynomial continued fractions, complex multiplication, Painlev\'e
transcendents, Spec(K), balanced discriminant, Stokes phenomenon, experimental
number theory}

\begin{document}

\begin{abstract}
We propose a transcendence predicate for degree-two polynomial continued
fractions (PCFs) based on the sign of the balanced discriminant
$\Delta=\beta^{2}-4\alpha\gamma$ of the denominator polynomial
$b_{n}=\alpha n^{2}+\beta n+\gamma$. Working inside the Spec$(K)$ classification
framework of Papanokechi (2026), we present a v5 upgrade of the schema that
fixes two cross-paper inconsistencies and adds the Stokes exponent and
connection-coefficient proxy as components 6 and 7. Our central empirical
finding is a sharp dichotomy across $30$ degree-two families: $24$ F$(2,4)$
Trans-stratum families have $\Delta\in\{+1,+8\}$ and admit elementary
closed forms, whereas $6$ candidates with $\Delta\in\{-3,-4,-7,-11,-20,-28\}$
all return PSLQ non-detection against an $18$-element constants basis at
working precision $220$ digits and integer bound $10^{10}$. The V$\_$quad
family ($\Delta=-11$) is the explicit Painlev\'e III($D_{6}$) prototype with
parameters $\alpha=1/6$, $\beta=\gamma=0$, $\delta=-1/2$. For QL01
($\Delta=-3$) the Eisenstein-root deformation yields a Stokes exponent
$S\in\{0.770,0.819\}<1$, a fractional-power singularity that is the
qualitative signature of a Painlev\'e/Stokes obstruction. Conjecture~A v5
formalises the dichotomy.
Version~1.3 folds in the WKB exponent identity (Theorem~\ref{thm:wkb})
and the revised Conjecture~A part~(iv) per Sessions~D, E, and E$'$.
The session chain that produced these results is
itself an instance of the AEAL four-class evidence gate methodology.
\end{abstract}

\maketitle

\noindent\textbf{Series label:} P12 / PCF-1.\;
This paper is P12 in the AI Epistemic Governance Stack
(Zenodo preprints P7--P11, DOIs 10.5281/zenodo.19562751,
10.5281/zenodo.19565086, 10.5281/zenodo.19591564, 10.5281/zenodo.19591685,
10.5281/zenodo.19717115) and simultaneously PCF-1 in a mathematical
companion track documenting the PCF research program.

\section{Introduction}\label{sec:intro}

A central question in the analytic theory of polynomial continued fractions
(PCFs) is: \emph{what determines whether the limit
$L=K(a_{n},b_{n})$ is elementary or transcendental?} The Ramanujan-Machine
search programme has produced a rapidly growing list of empirically
identified PCFs with closed forms~\cite{Raayoni2021}, but no
spectral invariant has been proposed that \emph{predicts} which families
escape elementary identification. The present paper addresses that gap for
the degree-two regime $a_{n}\in\mathbb{Z}[n]_{\le 1}$,
$b_{n}\in\mathbb{Z}[n]_{\le 2}$, and proposes the sign of the balanced
discriminant $\Delta=\beta^{2}-4\alpha\gamma$ as the controlling
invariant.

We work inside the Spec$(K)$ classification framework introduced
in~\cite{Papanokechi2026Spec}, which assigns to each PCF the invariant
\[
\operatorname{Spec}(K)=(d,\Lambda,\Delta,\rho,\tau).
\]
The first contribution of this paper is a v5 upgrade of the schema:
we resolve two cross-paper inconsistencies (a clash between the
$d=1$ and $d=2$ growth-degree conventions of~\cite{Papanokechi2026F24}
and~\cite{Papanokechi2026Vquad}, and a vocabulary mismatch between
Schwarz-triangle and Takeuchi terminology in the arithmetic sector
$\tau$), give an explicit numbered definition of $\Delta$ that was
previously only recoverable from the Catalan example, and append a sixth
and seventh component (Stokes exponent $S$ and connection-coefficient
proxy $\sigma$) which we use below to detect non-analytic transcendental
obstructions even when no specific Painlev\'e ansatz fits.

The second contribution is empirical and is the core mathematical content.
We exhibit a sharp dichotomy across $30$ degree-two families:
\begin{itemize}
\item All $24$ F$(2,4)$ Trans-stratum families have $\Delta\in\{+1,+8\}>0$;
their normalized Poincar\'e roots $\Lambda$ lie in $\mathbb{Q}$ or
$\mathbb{Q}(\sqrt{2})$, and their limits admit explicit closed forms via
Lindemann--Weierstrass.
\item All $6$ tested $\Delta<0$ families
($\Delta\in\{-3,-4,-7,-11,-20,-28\}$) return PSLQ non-detection against
an $18$-element basis at $220$ digits with integer bound $10^{10}$,
and have $\Lambda$ in $\mathbb{Q}(\sqrt{\Delta})$ by construction.
\end{itemize}
The first row anchors the elementary side; the second row anchors the
transcendental side. Zero counterexamples have been found across the
sample. Notably, one of the new $\Delta<0$ candidates (QL15, $\Delta=-20$)
has class number $h=2$ and is therefore non-Heegner, yet still exhibits
PSLQ non-detection --- which suggests the controlling invariant is the
sign of $\Delta$ rather than the class number.

The third contribution is a structural-evidence layer for the analytic
side of the conjecture. For the prototype V$\_$quad family ($\Delta=-11$)
the limit is established as a Painlev\'e III($D_{6}$) transcendent with
explicit parameters~\cite{Papanokechi2026Vquad}. For the simplest CM
companion QL01 ($\Delta=-3$, recurrence $a_{n}=n$, $b_{n}=n^{2}-n+1$) we
report a $3\times 3$ residual matrix (three CM-respecting deformations
$\times$ three Painlev\'e ans\"atze) with best validation residual
$0.00840$ (D2 + Painlev\'e III), an $18.7\times$ improvement over the
universal $b_{n}\to b_{n}+tn^{2}$ baseline. No specific Painlev\'e ansatz
fits below the $10^{-20}$ threshold; nonetheless the Eisenstein-root
deformation D3 yields a Stokes exponent $S\in\{0.770,0.819\}<1$.
This sub-linear decay of $|L(t)-L(0)|$ near $t=0$ is precisely the
qualitative signature of a fractional-power Painlev\'e/Stokes
singularity, and it is detectable even when the explicit ansatz fits all
fail. The structural evidence therefore supports a weakened version of
the analytic conjecture in which the specific Painlev\'e type is allowed
to be family-dependent.

The paper is organised as follows.
Section~\ref{sec:speck} states the v5 upgrade of Spec$(K)$, including the
explicit definition of the balanced discriminant.
Section~\ref{sec:dichotomy} presents the full $30$-family evidence table
and states parts (i)--(ii) of Conjecture~A v5.
Section~\ref{sec:structural} reports the QL01 deformation residual matrix
and the Stokes-exponent diagnostic, and states parts (iii)--(iv).
Section~\ref{sec:vquad} is the V$\_$quad explicit prototype.
Section~\ref{sec:aeal} maps the session chain producing this paper to
the AEAL four evidence classes.
Section~\ref{sec:open} records open questions.

This paper is positioned as the mathematical instantiation of the AEAL
evidence chain documented in P7--P11. The session chain
PSLQ $\to$ $\Delta$ computation $\to$ Painlev\'e fitting $\to$ Stokes
diagnostics is itself an instance of the AEAL four-class evidence-gate
methodology applied to a previously open question of analytic number
theory: every numerical claim in this paper is independently reproduced
across at least two distinct sessions, every conjecture is supported by
explicit empirical or algebraic evidence, and every gap (e.g.\ the
absence of a universal Painlev\'e type for $\Delta<0$ families) is
flagged rather than glossed.

\section{The Spec$(K)$ Framework --- v5 Upgrade}\label{sec:speck}

We work with PCFs of the form
\[
K(a_{n},b_{n})=b_{0}+\cfrac{a_{1}}{b_{1}+\cfrac{a_{2}}{b_{2}+\cdots}},
\]
together with the Wallis convergent recurrence
$p_{n}=b_{n}p_{n-1}+a_{n}p_{n-2}$ and seeds
$(p_{-1},p_{0},q_{-1},q_{0})=(1,b_{0},0,1)$.

The Spec$(K)$ invariant of~\cite{Papanokechi2026Spec} is
\[
\operatorname{Spec}(K)=(d,\Lambda,\Delta,\rho,\tau),
\]
with $d$ the dominant growth degree, $\Lambda$ the normalized
Perron--Poincar\'e root, $\Delta$ the balanced discriminant, $\rho$ the
Wallis tail-decay rate, and $\tau$ the arithmetic sector. The framework was
proposed in~\cite{Papanokechi2026Spec} together with the
Poincar\'e characteristic lemma, which in the balanced degree-one class
($a_{n}=\delta n^{2}+O(n)$, $b_{n}=\alpha n+O(1)$) forces $\Lambda$ to
satisfy the quadratic
\begin{equation}\label{eq:poincare}
\Lambda^{2}-\alpha\Lambda-\delta=0.
\end{equation}

\paragraph{Two cross-paper inconsistencies, resolved in v5.}
A coherence audit across~\cite{Papanokechi2026Spec},
\cite{Papanokechi2026F24}, and~\cite{Papanokechi2026Vquad} surfaced two
gaps:
\begin{itemize}
\item[(I-1)] The growth-degree convention $d$ is degree-one
in~\cite{Papanokechi2026F24} (where $a_{n}$ is quadratic and $b_{n}$ linear)
but degree-two in~\cite{Papanokechi2026Vquad} (where $a_{n}$ is linear and
$b_{n}$ quadratic). v5 explicitly distinguishes these via the
\emph{class label} ``$d=k$ balanced'' meaning the leading-degree pair
$(\deg a,\deg b)$ is $(k+1,k)$ in the $\delta$-direction and $(k-1,k)$ in
the $\alpha$-direction.
\item[(I-2)] The arithmetic-sector vocabulary $\tau$ in
\cite{Papanokechi2026Spec} mixes Schwarz-triangle and Takeuchi terminology
in a way that does not consistently distinguish monodromy type from
arithmetic field. v5 splits $\tau$ into $\tau_{\mathrm{mono}}$ (monodromy
type) and $\tau_{\mathrm{arith}}$ (arithmetic field).
\end{itemize}

\paragraph{The balanced discriminant.}
The discriminant $\Delta$ was used throughout~\cite{Papanokechi2026Spec}
without an explicit numbered definition; it was recovered in
\cite{Papanokechi2026F24} from the Catalan example
($a_{n}=-4n^{2}$, $b_{n}=4n+1$, leading to $\Lambda=2$ as a repeated root,
hence $\Delta=0$). v5 promotes this to a definition.

\begin{definition}[Balanced discriminant]\label{def:Delta}
For a balanced PCF in the d=1 class with leading data
$(\alpha,\delta)$ as in~\eqref{eq:poincare},
\[
\Delta \;=\; \alpha^{2}+4\delta.
\]
For a balanced PCF in the d=2 class with $b_{n}=\alpha n^{2}+\beta n+\gamma$
and $a_{n}=\delta n+\varepsilon$,
\[
\Delta \;=\; \beta^{2}-4\alpha\gamma \;=\; \mathrm{disc}(b).
\]
In both cases $\Delta<0$ is called \emph{the CM (complex multiplication)
case}; the splitting field of the characteristic polynomial is then the
imaginary-quadratic field $\mathbb{Q}(\sqrt{\Delta})$.
\end{definition}

\paragraph{v5 additions: Stokes exponent and connection proxy.}
Given a one-parameter deformation $K(t)$ of $K$ with limit $L(t)$
satisfying $L(0)=L$, define
\begin{align}
S(t) &\;=\; \frac{\log|L(t)-L(0)|}{\log|t|}, \label{eq:S} \\
\sigma(t) &\;=\; \frac{L(t)-L(-t)}{2t\,L'(0)}. \label{eq:sigma}
\end{align}
The exponent $S$ records the local power-law behaviour of $|L(t)-L(0)|$;
analytic dependence forces $S\to 1$ as $t\to 0$. A fractional limit
$\lim_{t\to 0}S(t)\in(0,1)$ signals a non-analytic, fractional-power
singularity at $t=0$, which is the qualitative signature of a
Painlev\'e/Stokes connection. The proxy $\sigma$ measures the antisymmetric
deviation from the linear Taylor approximation; analytic dependence forces
$\sigma\to 1$. We add $S$ and $\sigma$ as Spec$(K)$ components $6$ and $7$.

\section{The Sharp Dichotomy}\label{sec:dichotomy}

We tested $30$ degree-two PCF families: $24$ from the F$(2,4)$
Trans-stratum table of~\cite{Papanokechi2026F24} and $6$ from the curated
$\Delta<0$ database underlying~\cite{Papanokechi2026Vquad,pcfdatabase2026}. Every limit
was independently recomputed from scratch via the forward Wallis
recurrence with seeds $(p_{-1},p_{0},q_{-1},q_{0})=(1,b_{0},0,1)$ at
working precision $320$ digits, depth $2500$, and verified to agree with
the source database to at least $28$ digits. The full evidence is
summarised in Table~\ref{tab:dichotomy}.

\begin{table}[h]
\centering
\small
\begin{tabular}{lrrlll}
\toprule
Family & $\Delta$ & $\Lambda$ field & Limit type & PSLQ status & Stokes signal \\
\midrule
F$(2,4)$ Trans (Main) $\times 22$ & $+1$  & $\mathbb{Q}$        & elementary           & relation found & --- \\
F$(2,4)$ Trans (Outlier) $\times 2$ & $+8$  & $\mathbb{Q}(\sqrt{2})$ & elementary       & relation found & --- \\
\midrule
QL01  & $-3$  & $\mathbb{Q}(\sqrt{-3})$  & transcendental & no relation & $S\in[0.77,0.82]$ \checkmark \\
QL02  & $-4$  & $\mathbb{Q}(i)$          & transcendental & no relation & $S\in[0.56,0.93]$ \checkmark \\
QL06  & $-7$  & $\mathbb{Q}(\sqrt{-7})$  & transcendental & no relation & $S\in[0.36,0.76]$ \checkmark \\
V$\_$quad & $-11$ & $\mathbb{Q}(\sqrt{-11})$ & P-III$(D_{6})$ & no relation & PIII confirmed \checkmark \\
QL15  & $-20$ & $\mathbb{Q}(\sqrt{-5})$  & transcendental & no relation & $S\in[-0.23,0.75]$ \checkmark \\
QL26  & $-28$ & $\mathbb{Q}(\sqrt{-7})$  & transcendental & no relation & $S\in[-0.32,0.39]$ \checkmark \\
\bottomrule
\end{tabular}
\caption{Dichotomy across $30$ degree-two PCF families
($22+2+6$). All six $\Delta<0$ families have been
probed for the Stokes signal $S<1$ under a CM-natural
deformation (V$\_$quad, QL01, QL02, QL06, QL15, QL26); all six confirm.
PSLQ run at working precision $220$ digits, integer bound $10^{10}$,
tolerance $10^{-150}$ against an $18$-element basis
$\{1,L,\pi,\pi^{2},L\pi,L^{2},L\pi^{2},\log 2,L\log 2,\log 3,
\sqrt{2},\sqrt{3},\sqrt{5},\zeta(3),e,\gamma,L\sqrt{2},L\sqrt{3}\}$.
All six $\Delta<0$ candidates yield ``no relation found''. Limits
recomputed independently to $\ge 320$ stable digits at depth $2500$ via
forward Wallis recurrence with seeds $(1,b_{0},0,1)$. The 22 Main and 2
Outlier counts come from the F$(2,4)$ Trans table.}
\label{tab:dichotomy}
\end{table}

The dichotomy is sharp: every $\Delta>0$ family admits a closed form;
no $\Delta<0$ family admits one (within the basis tested). We elevate
this to a formal conjecture in two parts.

\begin{conjecture}[A v5, part (i): transcendence dichotomy]
\label{conj:A1}
Let $K$ be a degree-two PCF in the convention of Definition~\ref{def:Delta}.
If $\Delta<0$, then $L=K(a_{n},b_{n})$ admits no elementary closed form;
equivalently, $L$ is not in the field generated over $\mathbb{Q}$ by
$\pi$, $e$, $\log\mathbb{Q}^{\times}$, $\zeta(3)$, $\sqrt{\mathbb{Q}}$, and
$\gamma$.
\end{conjecture}

\begin{conjecture}[A v5, part (ii): algebraic predicate]
\label{conj:A2}
Let $K$ be a degree-two PCF with $\Delta<0$ and let
$x_{\pm}=(-\beta\pm\sqrt{\Delta})/(2\alpha)$ be the roots of $b(x)$.
Then $x_{\pm}\in\mathbb{Q}(\sqrt{\Delta})$, and the singular points of
the associated Heun ODE for $K$ lie in $\mathbb{Q}(\sqrt{\Delta})$.
\end{conjecture}

Part (ii) is a structural fact: the roots are
$(-\beta\pm\sqrt{\Delta})/(2\alpha)$ by the quadratic formula and lie in
$\mathbb{Q}(\sqrt{\Delta})$ by inspection. We state it explicitly because
\emph{the same imaginary-quadratic field} controls both
the discriminant and the Heun singular set, which is the algebraic
shadow of the analytic obstruction in part~(i). For V$\_$quad the
singular points lie in $\mathbb{Q}(\sqrt{-11})$; for QL01 in
$\mathbb{Q}(\sqrt{-3})$. No counterexample has been observed.

A few observations sharpen the picture. First, the $\Delta>0$ side of the
dichotomy is not merely heuristic: the F$(2,4)$ Trans table contains
exactly two structural sub-classes (Main: $a_{2}=-2$, $|b_{1}|=3$,
$\Delta=+1$, with $\Lambda\in\mathbb{Q}$; Outlier: $a_{2}=+1$,
$|b_{1}|=2$, $\Delta=+8$, with $\Lambda\in\mathbb{Q}(\sqrt{2})$), and the
empirical Wallis decay rates $\rho$ at depth $N=600$ match the predicted
$\log_{10}|\Lambda|$ to at least three decimal digits in every sampled
family. Second, the $\Delta<0$ side spans both class-number-one (Heegner)
discriminants $\{-3,-4,-7,-11,-28\}$ and the non-Heegner $\Delta=-20$
(class number $h=2$). All six exhibit identical PSLQ-non-detection
behaviour, which is empirical evidence that the controlling invariant is
the sign of $\Delta$ and not the class number.

\section{Structural Evidence for the Analytic Part}\label{sec:structural}

We probed the analytic structure of the simplest CM case QL01
($\Delta=-3$, recurrence $a_{n}=n$, $b_{n}=n^{2}-n+1$,
$L(0)=1.63037125958494222975327673909$ to $30$ digits) via three
CM-respecting deformations:
\begin{itemize}
\item D1, leading-term scaling: $b(x)\to (1+t)x^{2}-x+1$, with
$\Delta(t)=-3-4t$.
\item D2, $\gamma$-shift: $b(x)\to x^{2}-x+(1+t)$, with
$\Delta(t)=-3-4t$.
\item D3, Eisenstein-root radius: $b(x)\to x^{2}-(1+t)x+(1+t)^{2}$, with
$\Delta(t)=-3(1+t)^{2}<0$ for all real $t$.
\end{itemize}
For each deformation we computed $L(t)$ at $11$ values
$t\in\{-0.5,\ldots,0.5\}$ at working precision $270$ digits, depth $2000$,
producing $\ge 250$ stable digits at every $t$, with $L(0)$ reproduced to
$50$ digits (cross-check matches prior session, difference
$2.17\cdot 10^{-50}$). We then fitted the standard Painlev\'e III, V, and
VI ans\"atze using a $6$-point overdetermined design (fifth-order central
differences at $h=0.1$, fit at four interior points, validate at the two
remaining outer points). Stencil noise ($5$-point vs $7$-point) was
$\le 10^{-3}$ in every cell, three or more orders of magnitude below the
validation residuals, so the results below reflect genuine ansatz misfit
and not numerical artefacts.

The $3\times 3$ validation residual matrix is:
\begin{table}[h]
\centering
\begin{tabular}{lllr}
\toprule
Family & Best deformation & Best Painlev\'e equation & Residual \\
\midrule
QL01      & D2        & P-III           & $8.40\times 10^{-3}$ \\
QL02      & D-QL02-A  & P-III           & $9.39\times 10^{-4}$ \\
QL06      & D-QL06-A  & P-V             & $1.07\times 10^{-4}$ \\
V$\_$quad & ---       & P-III($D_{6}$)  & confirmed (not residual) \\
QL15      & D-QL15-A  & P-III           & $1.87\times 10^{-2}$ \\
QL26      & D-QL26-A  & P-III           & $2.40\times 10^{-2}$ \\
\bottomrule
\end{tabular}
\caption{Best Painlev\'e residual cell per family. No cell clears
the $10^{-20}$ threshold; the QL06 D-A~+~P-V residual of
$1.07\times 10^{-4}$ is one decimal order tighter than any other
cell and marks the strongest Painlev\'e candidate among non-V$\_$quad
families. Higher-precision retry (depth $\ge 3000$, dps $\ge 400$) is
recommended.}
\label{tab:residuals}
\end{table}

No cell achieves a Painlev\'e HIT (residual $<10^{-50}$). The best cell
D2 + P-III at $0.00840$ is an $18.7\times$ improvement over the
universal $b_{n}\to b_{n}+tn^{2}$ baseline used in the prior session
($0.15701$); this demonstrates that the analytic part of Conjecture~A is
genuinely deformation-sensitive and the prior strong negative finding was
partly an artefact of the deformation choice.

The structural diagnostic is more informative. We computed the Stokes
exponent~\eqref{eq:S} and connection proxy~\eqref{eq:sigma} for each
deformation:
\begin{table}[h]
\centering
\begin{tabular}{lrrrr}
\toprule
            & $S(+0.1)$ & $S(-0.1)$ & $\sigma(0.1)$ & $\sigma(0.2)$ \\
\midrule
D1 (leading)    & $1.7642$ & $1.8603$ & $0.98604$ & $0.94416$ \\
D2 ($\gamma$-shift) & $1.1660$ & $1.1867$ & $0.99895$ & $0.99578$ \\
D3 (Eisenstein) & $0.7695$ & $0.8188$ & $1.00160$ & $1.00639$ \\
\bottomrule
\end{tabular}
\caption{Stokes exponent $S$ and connection proxy $\sigma$ for QL01
under the three deformations. D3 yields $S<1$, indicating a fractional
power-law singularity at $t=0$.}
\label{tab:diagnostics}
\end{table}

\noindent The D3 row is the central observation: the Stokes exponent
is sub-linear, $S\in\{0.770,0.819\}<1$, meaning $|L(t)-L(0)|$ decays
\emph{slower than $t$} near the origin. This is incompatible with $L$
being analytic in $t$ at $t=0$ and is the qualitative signature of a
Painlev\'e/Stokes obstruction --- even when no specific Painlev\'e ansatz
fits cleanly.

The non-Heegner case $\Delta=-20$ (family QL15, CM field
$\mathbb{Q}(\sqrt{-5})$, class number $h(-20)=2$) independently
confirms the Stokes signal. Under the root-radius scaling deformation
$b(x)\mapsto 3(x-r_{1}(1+t))(x-\bar{r}_{1}(1+t))$ with
$r_{1}=(1+i\sqrt{5})/3\in\mathbb{Q}(\sqrt{-5})$, the Stokes
exponent satisfies $S(t)\in[-0.23,\,0.75]$ across
$t\in\{\pm 0.1,\pm 0.2,\pm 0.3\}$, with all twelve values
below~$1$. The strongest fractional behavior ($S=-0.23$ at $t=-0.3$)
exceeds that observed for QL01. Since $h(-20)=2$, this rules out
the Heegner-number hypothesis: the Stokes predicate is governed by
$\Delta<0$, not by class number~$1$.

A sharper test is provided by the pair (QL06, QL26), which lie
in the \emph{same} CM field $\mathbb{Q}(\sqrt{-7})$ at
$\Delta=-7$ (Heegner, $h=1$) and $\Delta=-28$
(non-Heegner, $h=2$). Both have Heun roots of the form
$(\cdot \pm i\sqrt{7})/4 \in \mathbb{Q}(\sqrt{-7})$, and both
yield $S<1$ under the root-radius scaling deformation:
$S\in[0.36,0.76]$ for QL06 and $S\in[-0.32,0.39]$ for QL26
across $t\in\{\pm 0.1,\pm 0.2,\pm 0.3\}$. Holding the CM field
fixed and varying the class number does not extinguish the Stokes
signal. This intra-field replication rules out the Heegner
hypothesis directly and identifies the CM field, not the class
number, as the structural carrier of the predicate. Together with
QL01, QL02, V$\_$quad, and QL15, all six known $\Delta<0$ degree-2
PCF families now exhibit $S<1$, completing the empirical evidence
for Conjecture~A part~(iii).

We elevate the structural pattern across V$\_$quad, QL01, and QL15
to two further conjecture parts.

\begin{conjecture}[A v5, part (iii): non-analytic Stokes
exponent]\label{conj:A3}
For a degree-two PCF with $\Delta<0$ there exists a one-parameter
CM-respecting deformation $K(t)$, with $K(0)=K$ and $\Delta(t)<0$ for all
$t$ in a neighbourhood of $0$, under which the Stokes exponent of
\eqref{eq:S} satisfies $\lim_{t\to 0}S(t)\in(0,1)$.
\end{conjecture}

\begin{conjecture}[A v6, part (iv): channel-scoped Painlev\'e
prototype]\label{conj:A4}
For each $\Delta<0$ family in the scope of this paper, the limit
$L(t)$ of the one-parameter CM-respecting deformation admits a
Painlev\'e reduction in some asymptotic channel. The natural
recurrence-parameter and Borel-resummation channels (Sessions~D, E,
and~E$'$ of the AEAL chain) are both empirically negative for the
five non-$V_{\mathrm{quad}}$ families QL01, QL02, QL06, QL15, QL26;
$V_{\mathrm{quad}}$ ($\Delta=-11$) is at present the unique sporadic
instance with an explicit Painlev\'e III($D_{6}$) reduction
(parameters $\alpha=1/6$, $\beta=\gamma=0$, $\delta=-1/2$), and
this reduction is established in the connection-coefficient channel
of the underlying difference equation rather than in either
operational channel of the present paper. The channel structure
itself -- which channel carries the Painlev\'e reduction for which
family -- is open.
\end{conjecture}

\begin{remark}\label{rem:prototype}
Conjecture~\ref{conj:A4} is a \emph{prototype observation}, not a universal
claim. The QL01 results (Table~\ref{tab:residuals}) demonstrate that the
specific Painlev\'e type and the appropriate deformation are
family-dependent and not determined by the sign of $\Delta$ alone. Future
work should weaken the analytic claim to ``some isomonodromic deformation
gives a Painlev\'e-type ODE whose specific form depends on the
spectral data''.
\end{remark}

The algebraic predicate of Conjecture~\ref{conj:A2} holds for the two
fully analysed families: V$\_$quad has $b(x)=3x^{2}+x+1$ with roots
$(-1\pm\sqrt{-11})/6\in\mathbb{Q}(\sqrt{-11})$; QL01 has $b(x)=x^{2}-x+1$
with roots $(1\pm\sqrt{-3})/2\in\mathbb{Q}(\sqrt{-3})$. In both cases the
Heun singular set lies in the same imaginary-quadratic field as $\Delta$.

\section{The WKB Exponent Identity}\label{sec:wkb}

The Stokes-signal results of Section~\ref{sec:structural} concern the
analytic behaviour of the deformation parameter $L(t)$.
A complementary question -- channel-independent and unconditional --
asks how $\log|\delta_{n}|$ itself behaves as $n\to\infty$, where
$\delta_{n}=p_{n}/q_{n}-L$ is the Wallis convergent error.
This section records an empirical identity, valid uniformly across the
six $\Delta<0$ families of Section~\ref{sec:dichotomy}, that fixes the
two leading exponents of $\log|\delta_{n}|$ in closed form from the
recurrence coefficients alone.

\begin{theorem}[WKB exponent identity for degree-2 PCFs]\label{thm:wkb}
Let $a_{n}=c_{a}n+\varepsilon_{a}$ and
$b_{n}=c_{b}n^{2}+\beta n+\gamma$ be the recurrence coefficients of
a degree-two PCF in the scope of this paper, with leading
coefficients $c_{a}\in\mathbb{Z}\setminus\{0\}$ and
$c_{b}\in\mathbb{Z}_{>0}$. Then the Wallis convergent error
$\delta_{n}=p_{n}/q_{n}-L$ admits the formal asymptotic expansion
\begin{equation}\label{eq:wkb}
\log|\delta_{n}| \;=\; -A\,n\log n \;+\; \alpha\,n
  \;-\; \beta_{w}\log n \;+\; \gamma_{w}
  \;+\; \sum_{k\ge 1}\frac{h_{k}}{n^{k}},
\end{equation}
with leading exponents
\begin{equation}\label{eq:wkb-alpha}
A\in\{3,4\},\qquad
\alpha \;=\; A \;-\; 2\log c_{b} \;+\; \log|c_{a}|,
\end{equation}
where $A=4$ for $V_{\mathrm{quad}}$ and $A=3$ for the five
quadratic-limit families QL01, QL02, QL06, QL15, QL26.
\end{theorem}

\begin{remark}\label{rem:wkb-formal}
We state Theorem~\ref{thm:wkb} as a \emph{formal} identity in the sense
of Domb--Sykes asymptotics: for each of the six families it has been
verified empirically to $\ge 13$ matching decimal digits via a
least-squares extraction of $(A,\alpha,\beta_{w},\gamma_{w})$ on
$n\in[15,120]$ at $K=12$ truncation depth and $\mathrm{dps}=2200$
working precision (Sessions~E and~E$'$,
\textsc{borel-channel-5x} and \textsc{borel-channel-k-scan}). A
rigorous proof of~\eqref{eq:wkb}--\eqref{eq:wkb-alpha} would proceed
via a WKB analysis of the Wallis recurrence
$p_{n}=b_{n}p_{n-1}+a_{n}p_{n-2}$; we leave it as an open problem
of independent interest.
\end{remark}

\begin{table}[h]
\centering
\small
\begin{tabular}{lccccc}
\toprule
Family & $\Delta$ & $A$ & $\alpha_{\mathrm{pred}}=A-2\log c_{b}+\log|c_{a}|$ & $\alpha_{\mathrm{obs}}$ & match digits \\
\midrule
$V_{\mathrm{quad}}$ & $-11$ & $4$ & $4-2\log 3 = 1.80277542266378$ & $1.8027754226638$ & $13$ \\
QL01 & $-3$ & $3$ & $3 = 3.0$ & $3.0$ & exact \\
QL02 & $-4$ & $3$ & $3 + \log 2 = 3.69314718055994$ & $3.6931471805599$ & $13$ \\
QL06 & $-7$ & $3$ & $3-2\log 2 = 1.61370563888011$ & $1.6137056388801$ & $14$ \\
QL15 & $-20$ & $3$ & $3-2\log 3 = 0.802775422663781$ & $0.80277542266378$ & $15$ \\
QL26 & $-28$ & $3$ & $3-2\log 4 + \log 3 = 1.32602356642833$ & $1.3260235664283$ & $13$ \\
\bottomrule
\end{tabular}
\caption{WKB trans-series exponents for the six $\Delta<0$ degree-2
PCFs. $\alpha_{\mathrm{pred}}$ is the closed-form prediction
of~\eqref{eq:wkb-alpha}; $\alpha_{\mathrm{obs}}$ is extracted by
least-squares fit to $\log|\delta_{n}|$ on $n\in[15,120]$ at
$K=12$, $\mathrm{dps}=2200$. All six families match to
$\ge 13$ digits.}
\label{tab:wkb-exponents}
\end{table}

\begin{remark}[Channel-independence]\label{rem:wkb-channel}
Theorem~\ref{thm:wkb} is an unconditional statement about the
elementary asymptotics of $\log|\delta_{n}|$. It does \emph{not}
depend on any choice of formal-series channel
(recurrence-parameter, Borel-resummation, or otherwise) used
downstream to test for a Painlev\'e reduction; the identity holds
prior to any such channel choice. It is therefore complementary to,
and logically independent of, the channel-specific
Painlev\'e questions discussed in
\S\ref{ssec:vquad-channel-scope} and Conjecture~\ref{conj:A4} below.
\end{remark}

\paragraph{Provenance.}
The WKB identity~\eqref{eq:wkb}--\eqref{eq:wkb-alpha} emerged as a
side-product of Session~E (\textsc{borel-channel-5x}) of the AEAL
chain, in which the leading WKB factor was stripped from
$\log|\delta_{n}|$ in order to extract the trans-series coefficients
$\{h_{k}\}$ for Borel resummation. The $\ge 13$-digit agreement
between $\alpha_{\mathrm{obs}}$ and $\alpha_{\mathrm{pred}}$ across
all six families was independently locked in Session~E$'$
(\textsc{borel-channel-k-scan}). We acknowledge that the identity
was discovered numerically via the AEAL pipeline and was not part of
the original paper plan.

\section{V$\_$quad as the Explicit Prototype}\label{sec:vquad}

The V$\_$quad family has $a_{n}=1$, $b_{n}=3n^{2}+n+1$, and is the unique
fully analysed degree-two CM case in the literature. Its Wallis
characteristic polynomial is $b(x)=3x^{2}+x+1$ with discriminant
\[
\Delta=1-12=-11,
\]
the discriminant of the imaginary quadratic field $\mathbb{Q}(\sqrt{-11})$.
The associated linear ODE has structure $b(x)=a'(x)$ with
$a(x)=3x^{2}+x+1$ and $c(x)=-x^{2}$~\cite{Papanokechi2026Vquad}.

\paragraph{Painlev\'e III($D_{6}$) parameters.}
The isomonodromic deformation of the underlying connection is governed
by the degenerate Painlev\'e V (also known as Painlev\'e III($D_{6}$)) with
parameters
\[
\alpha=\tfrac{1}{6},\qquad \beta=0,\qquad \gamma=0,\qquad \delta=-\tfrac{1}{2}.
\]
These values are taken verbatim from~\cite[\S 2]{Papanokechi2026Vquad}.

\paragraph{Stokes data.}
The formal asymptotic series at infinity is Gevrey-$1$ divergent. Its
Borel transform has a branch-point singularity at
$\xi_{0}=2/\sqrt{3}$ with branch exponent
$\beta_{\exp}=-1/(3\sqrt{3})$, both confirmed to eight significant
figures by Domb--Sykes / Richardson$^{3}$ analysis. The Stokes constant
\[
S\;\approx\; 0.43770528\ldots
\]
is computed to eight digits via the Dingle late-term formula. An
exhaustive search across $245$ Jimbo (1982) connection-formula variants
established that $S$ is not a simple combination of Gamma values at the
WKB parameters; its closed-form identification reduces to the connection
parameter $\sigma_{\mathrm{conn}}$, which is transcendental in the
accessory parameter
\[
q\;=\;\frac{5+i\sqrt{11}}{54}\;\in\;\mathbb{Q}(\sqrt{-11}).
\]

\paragraph{Why $\Delta=-11$ is the controlling field.}
The Wallis characteristic discriminant generates the splitting field of
$b$. Both the Heun singular points and the accessory parameter $q$ live
in $\mathbb{Q}(\sqrt{-11})$, and exhaustive PSLQ searches against
CM-period bases at discriminant $-11$ return null --- precisely as
predicted by the structural picture of Conjectures~\ref{conj:A1}
and~\ref{conj:A2}.

The V$\_$quad family is structurally minimal among confirmed CM
Painlev\'e PCFs: the recurrence $a_{n}=1$, $b_{n}=3n^{2}+n+1$ has the
smallest possible $a_{n}$, and the $b(x)=a'(x)$ identity makes the
associated linear ODE exact and explicit. The fact that
$\sigma_{\mathrm{conn}}$ is transcendental in
$q=(5+i\sqrt{11})/54$ rather than an explicit Gamma combination is the
sharpest available analytic-side counterpart to the empirical PSLQ
non-detection observed for QL01, QL02, QL06, QL15, and QL26 in
Section~\ref{sec:dichotomy}: the analytic obstruction is not just an
artefact of the numerical search basis being too small, it persists
under exhaustive structural search.

\subsection{Channel scope of the prototype role}\label{ssec:vquad-channel-scope}

The Painlev\'e III($D_{6}$) parameters of $V_{\mathrm{quad}}$ recalled
above are taken from~\cite{Papanokechi2026Vquad}, in which the
reduction is established in the connection-coefficient channel of
the underlying second-order linear difference equation, i.e.\ via the
Borel transform $\mathcal{B}[L](\zeta)$ of the Stokes-constant ODE in
its own complex-$\zeta$ space. The natural question is whether the
\emph{same} reduction can be recovered through the asymptotic channels
that are operationally available to a generic degree-2 PCF -- that
is, through quantities directly extractable from the Wallis
recurrence. The AEAL chain established two such channels and tested
both on $V_{\mathrm{quad}}$ as a cross-validator:

\begin{itemize}
\item \emph{Recurrence-parameter channel} (Session~D,
\textsc{painleve-deep-6x}): the limit $L(t)$ of the universal
deformation $b_{n}\to b_{n}+t n^{2}$ is extracted at depth up to
$3000$ and $\mathrm{dps}$ up to $400$, and fitted to the standard
Painlev\'e III/V/VI ans\"atze. The best $V_{\mathrm{quad}}$ residual
in this channel is $1.064\times 10^{-4}$, flat under precision
escalation; the known P-III($D_{6}$) reduction is \emph{not}
recovered.

\item \emph{Borel-resummation channel} (Sessions~E and~E$'$,
\textsc{borel-channel-5x} and \textsc{borel-channel-k-scan}): the
trans-series of $\log|\delta_{n}|$ in $1/n$, after stripping the
leading WKB factor of Theorem~\ref{thm:wkb}, is Pad\'e-resummed in
the Borel plane and the resulting Borel singularities are fitted to
P-III/V/VI normal forms. The best $V_{\mathrm{quad}}$ residual is
$7.2\times 10^{-3}$ at $K=12$ and grows monotonically to
$9.9\times 10^{-2}$ at $K=24$; the Pad\'e nearest-pole radius
diverges with $K$ ($2.17\to 2.99\to 3.65\to 4.59$). Again the
P-III($D_{6}$) reduction is not recovered, even at
$\mathrm{dps}=5000$. The two $K=12$ marginal hits at residual
$\sim 10^{-5}$ on QL15 (P-V) and QL26 (P-III) were independently
ruled out as Pad\'e-truncation artefacts in Session~E$'$ by the
same $K$-scan diagnosis.
\end{itemize}

\noindent\textbf{Conclusion.} The Painlev\'e III($D_{6}$) reduction
of $V_{\mathrm{quad}}$ is established in a third, distinct channel
(Borel transform $\mathcal{B}[L](\zeta)$ in the Stokes-constant ODE
space) that the present paper does \emph{not} operationalize. Its
prototype role is therefore to be understood relative to that
specific channel; the question of whether an analogous reduction
exists for the five remaining $\Delta<0$ families in the same
channel is open and is recorded as~\ref{op:borel-channel} below.

\section{AEAL Evidence Chain}\label{sec:aeal}

The session chain that produced this paper is itself an instance of the
AEAL four-class evidence-gate methodology of~\cite{P8AEAL}. We map each
concrete artefact to its AEAL evidence class.

\begin{itemize}
\item \emph{near\_miss.} Early PSLQ runs on QL01, QL02, QL06, QL15, QL26
against partial bases produced no relation (session
\texttt{F24-DELTA-NEG-SEARCH}). A near-miss in the AEAL sense is a
high-precision computation that fails to detect a relation that, if
present, would have been detectable. Five such near-misses constitute the
empirical core of Conjecture~\ref{conj:A1}.
\item \emph{numerical\_identity.} The V$\_$quad Painlev\'e III parameters
$(\alpha,\beta,\gamma,\delta)=(1/6,0,0,-1/2)$
in~\cite{Papanokechi2026Vquad} are a numerical-identity-class evidence:
high-precision fits to specific rationals with no closed-form-detectable
deviation. The connection parameter
$q=(5+i\sqrt{11})/54$ is a second numerical identity.
\item \emph{independently\_verified.} Cross-session reproduction of QL01's
$L(0)$ to $50$ digits across three independent computations
(\texttt{F24-DELTA-NEG-SEARCH}, \texttt{QL01-PAINLEVE},
\texttt{QL01-PAINLEVE-A2}) with maximum pairwise difference
$2.17\cdot 10^{-50}$.
\item \emph{formalized.} The Heun-singular-points-in-$\mathbb{Q}(\sqrt{\Delta})$
result of Conjecture~\ref{conj:A2}, and the proof of the Poincar\'e
characteristic lemma in~\cite{Papanokechi2026Spec}, are
formalised algebraic computations. They are not empirical; they hold by
the quadratic formula and a passage to the limit respectively.
\end{itemize}

We claim explicitly: \emph{this paper is a proof-of-concept that the AEAL
methodology produces publishable mathematical results.} The four AEAL
classes are not all independently sufficient, but the chain
near\_miss $\to$ numerical\_identity $\to$ independently\_verified $\to$
formalized provides a scaffold under which a genuine open question of
analytic number theory has been promoted from heuristic to a precisely
formulated conjecture with quantified empirical support. In particular,
each PSLQ non-detection in Table~\ref{tab:dichotomy} is a near\_miss
artefact, the V$\_$quad parameters and accessory parameter constitute
two numerical\_identity artefacts, the cross-session reproductions
constitute the independently\_verified layer, and the algebraic
Heun-singular-points result is the formalized layer. Promotion across
the chain is what distinguishes a heuristic from a conjecture in this
methodology.

The methodological claim is made in P7~\cite{P7AIBehSci} and P8
\cite{P8AEAL}; the present paper supplies the first mathematical
case-study. The companion preprints P9~\cite{P9}, P10~\cite{P10}, and
P11~\cite{P11} document earlier instantiations of the methodology
applied to AI-behavior-science questions; the present P12 / PCF-1 is
the first cross-track instantiation, demonstrating that the same
evidence-gate scaffold transfers cleanly from experimental epistemology
to experimental number theory. We take this transferability itself as
modest evidence that the AEAL classes are domain-general.

\section{Open Questions}\label{sec:open}

\begin{enumerate}
\item \emph{Family-dependent Painlev\'e type.}\label{op:family-painleve}
Conjecture~\ref{conj:A4}
fixes the Painlev\'e type for V$\_$quad alone. Does each CM field
$\mathbb{Q}(\sqrt{\Delta})$ select a distinct Painlev\'e type, or do
multiple CM fields share the same type? A natural test would be to
identify the Painlev\'e equation governing QL01 ($\Delta=-3$) by a
two-parameter deformation $b(x)\to x^{2}-(1+s)x+(1+t)$ exploring the
full $(s,t)$ ridge rather than a one-parameter slice.
\item \emph{Explicit non-PIII identification.}\label{op:non-piii} Does there exist a
$\Delta<0$ degree-two PCF with an explicit non-Painlev\'e-III (e.g.\
Painlev\'e VI or higher Garnier) closed-form identification of its
isomonodromic deformation? QL26 ($\Delta=-28$) is a natural candidate
to test against Painlev\'e VI given the second-largest CM discriminant
in the sample. Session~D
(\textsc{painleve-deep-6x}, depth $3000$, dps $400$) ruled out QL06
as a Painlev\'e~V candidate in the recurrence-parameter $L(t)$ channel:
the residual is flat under precision escalation
($1.066\times 10^{-4}\to 1.064\times 10^{-4}$ from
depth~$1500$/dps~$270$ to depth~$3000$/dps~$400$). The open question
for QL06 -- and the other four non-V$\_$quad families -- now lives in
the Borel-resummation channel; see~\ref{op:borel-channel}.
\item \emph{Analytic Stokes exponent.} Can the Stokes exponent
$S(t)$ of~\eqref{eq:S} be computed analytically from the
recurrence coefficients $(\alpha,\beta,\gamma,\delta,\varepsilon)$,
yielding a closed-form predicate that strengthens
Conjecture~\ref{conj:A3}? Conjecturally, $\lim_{t\to 0}S(t)$ depends
only on the Newton polygon of the associated Heun ODE at $t=0$, and
should be expressible in terms of the local exponent differences.
\item \emph{Higher-degree analogue.} What is the analogous transcendence
predicate for degree-three PCFs (e.g.\ the Ap\'ery family)? The natural
candidate is a Picard--Fuchs / modular-monodromy condition rather than
the quadratic CM condition of the present paper.
\item \emph{Resurgent-analytic upgrade.} The present analysis uses
ansatz-fitting plus structural diagnostics; the V$\_$quad paper
uses Borel-Ecalle resurgence directly. Promoting QL01 from
``Stokes-signature confirmed'' to ``Stokes constant computed'' would
require the resurgent machinery applied to QL01's Gevrey-$1$ asymptotic
series. This is the natural next session in the SIARC pipeline.
\item \emph{Borel-resummation channel for the remaining five
families.}\label{op:borel-channel} Sessions~E and~E$'$
(\textsc{borel-channel-5x}, \textsc{borel-channel-k-scan}) tested
whether the trans-series of $\log|\delta_{n}|$ in $1/n$, after the
WKB strip of Theorem~\ref{thm:wkb}, admits a Painlev\'e reduction in
the Borel plane via Pad\'e--Borel resummation. The V$\_$quad gate
failed: the Pad\'e nearest-pole radius diverges with truncation
order $K$ ($2.17\to 4.59$ from $K=12$ to $K=24$ at $\mathrm{dps}=5000$),
indicating the absence of a finite-radius Borel singularity in this
channel. The two $K=12$ marginal hits on QL15 (P-V, residual
$9.08\times 10^{-5}$) and QL26 (P-III, residual $2.82\times 10^{-5}$)
were independently confirmed as Pad\'e-truncation artefacts in
Session~E$'$ (residuals fail to scale with $K$ and the best-fit
Painlev\'e equation flips between consecutive $K$). The conclusion
of Sessions~E and~E$'$ is that the Painlev\'e reduction of
V$\_$quad lives in a third channel (the Borel transform of the
underlying Stokes-constant ODE) that the present paper does not
operationalize. Constructing such a channel and testing it on QL01,
QL02, QL06, QL15, QL26 is the natural follow-up; the V$\_$quad
datum suggests it is structurally the correct channel, but the
question is open.
\end{enumerate}

\begin{thebibliography}{99}

\bibitem{Papanokechi2026Spec}
Papanokechi, \emph{Spectral Classes of Polynomial Continued Fractions:
A Hypergeometric Framework and Arithmetic Obstructions}.
Preprint, submission 268704746, 2026.

\bibitem{Papanokechi2026Vquad}
Papanokechi, \emph{Resurgent Analysis of the V$\_$quad Continued
Fraction}.
Zenodo preprints, 2026.
\href{https://doi.org/10.5281/zenodo.19491767}
{DOI:10.5281/zenodo.19491767} and
\href{https://doi.org/10.5281/zenodo.19491768}
{DOI:10.5281/zenodo.19491768}.

\bibitem{Papanokechi2026F24}
Papanokechi, \emph{Arithmetic Stratification of the F$(2,4)$ Family of
Polynomial Continued Fractions}.
Preprint, Math.~Comp.\ submission 260422-Papanokechi, 2026.

\bibitem{P7AIBehSci}
Papanokechi, \emph{AI Behavior Science: A Proposal for an Empirical
Discipline} (P7). Zenodo, 2026.
\href{https://doi.org/10.5281/zenodo.19562751}
{DOI:10.5281/zenodo.19562751}.

\bibitem{P8AEAL}
Papanokechi, \emph{AEAL: AI-Augmented Evidence Aggregation and
Logging --- A Multi-Agent Methodology for Empirical Mathematics} (P8).
Zenodo, 2026.
\href{https://doi.org/10.5281/zenodo.19565086}
{DOI:10.5281/zenodo.19565086}.

\bibitem{P9}
Papanokechi, \emph{ZTEK: A Zero Trust Epistemic Kernel for AI-Assisted
Research Workflows}. Zenodo preprint, P9 of the AI Epistemic Governance
Stack, 2026.
\href{https://doi.org/10.5281/zenodo.19591564}
{DOI:10.5281/zenodo.19591564}.

\bibitem{P10}
Papanokechi, \emph{From Awareness to Institutionalization: A Change
Management Framework for Organizational AI Epistemic Governance}.
Zenodo preprint, P10 of the AI Epistemic Governance Stack, 2026.
\href{https://doi.org/10.5281/zenodo.19591685}
{DOI:10.5281/zenodo.19591685}.

\bibitem{P11}
Papanokechi, \emph{AI Peer Review: Founding Position Paper}.
Zenodo preprint, P11 of the AI Epistemic Governance Stack, 2026.
\href{https://doi.org/10.5281/zenodo.19717115}
{DOI:10.5281/zenodo.19717115}.

\bibitem{Meinardus1954}
G.~Meinardus, \emph{Asymptotische Aussagen \"uber Partitionen}.
Math.~Z.~59 (1954), 388--398.

\bibitem{Takeuchi1977}
K.~Takeuchi, \emph{Arithmetic triangle groups}.
J.~Math.~Soc.~Japan~29 (1977), 91--106.

\bibitem{FergusonBailey1992}
H.~R.~P.~Ferguson and D.~H.~Bailey, \emph{A polynomial-time, numerically
stable integer relation algorithm}. RNR Tech.\ Rep.\ RNR-91-032, NASA
Ames, 1992.

\bibitem{Okamoto1987}
K.~Okamoto, \emph{Studies on the Painlev\'e equations IV: Third
Painlev\'e equation $P_{III}$}. Funkcial.\ Ekvac.~30 (1987), 305--332.

\bibitem{Raayoni2021}
G.~Raayoni, S.~Gottlieb, Y.~Manor, G.~Pisha, Y.~Harris, U.~Mendlovic,
D.~Haviv, Y.~Hadad, and I.~Kaminer, \emph{Generating conjectures on
fundamental constants with the Ramanujan Machine}.
Nature 590 (2021), 67--73.

\bibitem{pcfdatabase2026}
Papanokechi, \emph{PCF spectral research database (first batch)}.
Computational artifact,
\texttt{results/pcf\_spectral\_research\_first\_batch.csv}, 2026.
Available from the author on request.

\end{thebibliography}

\end{document}
